\subsection{Applications: The Welfare Function in Practice}
\label{subsec:applications}

%%%%%%%%%%%%%%%%%%%%%%%%%%%%%%%%%%%%%%%%%%%%%%%%%%%%%%%%%%%%%%%%%%%%%%%%%%%%%%%
% 10.9: Nine application domains demonstrating the welfare function
%%%%%%%%%%%%%%%%%%%%%%%%%%%%%%%%%%%%%%%%%%%%%%%%%%%%%%%%%%%%%%%%%%%%%%%%%%%%%%%

The previous sections developed the EBF welfare function theoretically.
This section demonstrates its practical application across nine domains,
each illustrating different aspects of the 144-component structure,
inter-category complementarities, and context modulation.

\subsubsection{Overview: Nine Application Domains}

\begin{table}[htbp]
\centering
\caption{Nine Application Domains}
\label{tab:nine-domains}
\renewcommand{\arraystretch}{1.3}
\begin{tabular}{@{}clll@{}}
\toprule
\textbf{\#} & \textbf{Domain} & \textbf{Core Theme} & \textbf{Primary Conflict} \\
\midrule
10.9.1 & Labor Market Transitions & Career change & INU↑ vs. IDN↓ \\
10.9.2 & Health \& Prevention & Lifestyle change & $t=0$ vs. $t=3$ \\
10.9.3 & Financial Decisions & Saving behavior & F now vs. F later \\
10.9.4 & Sustainability \& Climate & Pro-environmental & INU vs. KNU/Eco \\
10.9.5 & Digital Transformation & Organizational change & D↑, S?, IDN? \\
10.9.6 & Organizational Culture & Trust \& cooperation & KNU, $\gamma$-structure \\
10.9.7 & Education \& Human Capital & Lifelong learning & C, time trade-off \\
10.9.8 & Public Policy & State intervention & Macro application \\
10.9.9 & Marketing \& Consumer & Purchase decisions & Awareness manipulation \\
\bottomrule
\end{tabular}
\end{table}

Each domain section follows a consistent structure:
\begin{enumerate}
\item Context description ($\Psi$-profile)
\item Relevant utility components (which of the 144?)
\item Typical conflicts (INU vs. KNU vs. IDN)
\item Reference point dynamics
\item Complementarity structure
\item Intervention levers
\item Concrete case example
\end{enumerate}

%%%%%%%%%%%%%%%%%%%%%%%%%%%%%%%%%%%%%%%%%%%%%%%%%%%%%%%%%%%%%%%%%%%%%%%%%%%%%%%
\subsubsection{Domain 10.9.1: Labor Market Transitions}
\label{subsubsec:labor-transitions}
%%%%%%%%%%%%%%%%%%%%%%%%%%%%%%%%%%%%%%%%%%%%%%%%%%%%%%%%%%%%%%%%%%%%%%%%%%%%%%%

\paragraph{The Challenge.}
Economic transformation---automation, globalization, green transition---requires
workers to change occupations, industries, or skill sets. Despite financial
incentives, many workers resist transitions that would improve their economic
position. The EBF framework explains why.

\paragraph{Context Profile ($\Psi$).}

\begin{table}[htbp]
\centering
\caption{Labor Transition Context Factors}
\label{tab:labor-psi}
\renewcommand{\arraystretch}{1.3}
\begin{tabular}{@{}lll@{}}
\toprule
\textbf{$\Psi$-Dimension} & \textbf{Typical Level} & \textbf{Effect on Transition} \\
\midrule
$\Psi_I$ (Institutional) & Variable & High: safety nets enable risk-taking \\
$\Psi_S$ (Social trust) & Often LOW in declining regions & Low: community fragmentation \\
$\Psi_C$ (Cognitive) & Variable & High: easier skill acquisition \\
$\Psi_T$ (Temporal) & LOW during disruption & Low: uncertainty blocks planning \\
$\Psi_F$ (Flexibility) & Critical & High: cultural openness to change \\
\bottomrule
\end{tabular}
\end{table}

\paragraph{Relevant Utility Components.}
The transition decision activates specific components across all three categories:

\begin{table}[htbp]
\centering
\caption{Labor Transition: Key Utility Components}
\label{tab:labor-components}
\small
\renewcommand{\arraystretch}{1.2}
\begin{tabular}{@{}llll@{}}
\toprule
\textbf{Component} & \textbf{Old Job} & \textbf{New Job} & \textbf{$\Delta$} \\
\midrule
\multicolumn{4}{l}{\textit{Individual Utility (INU)}} \\
$U^{INU}_{F,2}$ (Financial, medium-term) & \euro{45k} & \euro{65k} & \textbf{+\euro{20k}} \\
$U^{INU}_{P,2}$ (Physical, medium-term) & 0.4 (hazardous) & 0.7 (safe) & \textbf{+0.3} \\
$U^{INU}_{E,1}$ (Emotional, short-term) & 0.5 & 0.3 (stress) & \textbf{--0.2} \\
\midrule
\multicolumn{4}{l}{\textit{Collective Utility (KNU)}} \\
$U^{KNU}_{S,0}$ (Social, instant) & 0.8 (community) & 0.3 (alone) & \textbf{--0.5} \\
$U^{KNU}_{F,2}$ (Family financial) & 0.5 & 0.7 & \textbf{+0.2} \\
\midrule
\multicolumn{4}{l}{\textit{Identity Utility (IDN)}} \\
$U^{IDN}_{S,0}$ (Social identity) & 0.9 (``I'm a miner'') & 0.2 & \textbf{--0.7} \\
$U^{IDN}_{P,0}$ (Physical identity) & 0.8 (``I work with hands'') & 0.3 & \textbf{--0.5} \\
$U^{IDN}_{E,2}$ (Meaning, medium-term) & 0.7 & 0.4 & \textbf{--0.3} \\
\bottomrule
\end{tabular}
\end{table}

\paragraph{The Core Conflict: INU↑ vs. IDN↓.}
The transition offers clear INU gains (money, safety) but severe IDN losses:

\begin{align}
\Delta U^{INU} &\approx +0.4 \quad \text{(substantial gain)} \\
\Delta U^{KNU} &\approx -0.2 \quad \text{(moderate loss)} \\
\Delta U^{IDN} &\approx -0.8 \quad \text{(severe loss)}
\end{align}

\paragraph{Reference Point Analysis.}

\begin{table}[htbp]
\centering
\caption{Reference Points in Labor Transition}
\label{tab:labor-reference}
\renewcommand{\arraystretch}{1.3}
\begin{tabular}{@{}llp{6cm}@{}}
\toprule
\textbf{Category} & \textbf{Reference Point} & \textbf{Formation} \\
\midrule
INU & $r^{INU}_F = $ \euro{45k} & Status quo salary; peers in same job \\
KNU & $r^{KNU}_S = 0.8$ & Current community cohesion level \\
IDN & $r^{IDN} = $ ``Coal miner'' & Decades of identity formation; family tradition; community role \\
\bottomrule
\end{tabular}
\end{table}

The IDN reference point is \emph{stickiest}---it doesn't adapt to the new 
reality because identity is tied to the old role.

\paragraph{Complementarity Structure.}
Before transition (status quo):
\begin{align}
\gamma_{INU,KNU} &> 0 \quad \text{(my work supports community)} \\
\gamma_{INU,IDN} &> 0 \quad \text{(my income confirms ``provider'' identity)} \\
\gamma_{KNU,IDN} &> 0 \quad \text{(``we are miners together'')}
\end{align}

After transition:
\begin{align}
\gamma_{INU,KNU} &< 0 \quad \text{(my success = leaving community)} \\
\gamma_{INU,IDN} &< 0 \quad \text{(financial gain = ``selling out'')} \\
\gamma_{KNU,IDN} &< 0 \quad \text{(betraying ``us'')}
\end{align}

\textbf{The transition inverts the entire complementarity structure!}

\paragraph{Total Welfare Calculation.}
Using typical weights ($\omega_{INU} = 0.4$, $\omega_{KNU} = 0.25$, $\omega_{IDN} = 0.35$):

\begin{align}
\Delta Q &= \omega_{INU} \cdot \Delta U^{INU} + \omega_{KNU} \cdot \Delta U^{KNU} + \omega_{IDN} \cdot \Delta U^{IDN} \nonumber \\
&\quad + \Delta(\gamma_{INU,KNU} \cdot U^{INU} \cdot U^{KNU}) + \ldots \nonumber \\
&= 0.4 \cdot (+0.4) + 0.25 \cdot (-0.2) + 0.35 \cdot (-0.8) + \text{complementarity collapse} \nonumber \\
&= 0.16 - 0.05 - 0.28 + (\approx -0.15) \nonumber \\
&= \mathbf{-0.32}
\end{align}

\textbf{Result:} Despite +\euro{20k} salary increase, total welfare \emph{decreases}.
The rational worker rejects the transition.

\paragraph{Intervention Strategies.}

\begin{table}[htbp]
\centering
\caption{Intervention Strategies for Labor Transitions}
\label{tab:labor-interventions}
\renewcommand{\arraystretch}{1.3}
\begin{tabular}{@{}llll@{}}
\toprule
\textbf{Lever} & \textbf{Target} & \textbf{Strategy} & \textbf{Example} \\
\midrule
Nudge (A) & $a_{INU}$ ↑ & Highlight financial benefits & ``Your family needs this'' \\
Educate ($\omega$) & $\omega_{IDN}$ ↓ & Reduce identity rigidity & Growth mindset training \\
Design ($\Delta U$) & $\Delta U^{IDN}$ ↑ & Create identity bridge & ``Green energy worker'' \\
Design ($\Delta U$) & $\Delta U^{KNU}$ ↑ & Collective transition & Move community together \\
Context ($\Psi$) & $\Psi_I$ ↑ & Safety nets & Retraining with income guarantee \\
\bottomrule
\end{tabular}
\end{table}

\paragraph{The Identity Bridge Solution.}
The most effective intervention addresses IDN directly:

\begin{quote}
\textit{``You're not becoming a coder. You're becoming a \textbf{green energy technician}---
someone who builds the future with their hands, just like you've always done.
The community needs people like you to lead this transition.''}
\end{quote}

This reframes the transition to:
\begin{itemize}
\item Preserve core identity elements (``hands-on work'', ``building things'')
\item Create new identity that honors the old (``leading the transition'')
\item Maintain KNU by making it a collective journey
\end{itemize}

\paragraph{Case Example: German Coal Transition.}
Germany's planned coal phase-out includes:
\begin{itemize}
\item \euro{40 billion structural fund for affected regions
\item Retraining programs framed as ``energy transition careers''
\item Community-level transition planning (preserving KNU)
\item Explicit identity work: ``From coal to climate''
\end{itemize}

Early evidence suggests higher acceptance than purely financial approaches.

\paragraph{Key Insight.}
\begin{tcolorbox}[colback=yellow!10!white, colframe=yellow!50!black]
\textbf{Labor Market Insight:}
Financial incentives alone fail because they address only INU.
Successful transitions require:
\begin{enumerate}
\item Identity bridges (IDN preservation)
\item Collective movement (KNU maintenance)
\item Complementarity preservation ($\gamma > 0$)
\end{enumerate}
The welfare function reveals why rational workers reject ``irrational'' 
choices---and how to design better policy.
\end{tcolorbox}

%%%%%%%%%%%%%%%%%%%%%%%%%%%%%%%%%%%%%%%%%%%%%%%%%%%%%%%%%%%%%%%%%%%%%%%%%%%%%%%
% END OF 10.9.1
%%%%%%%%%%%%%%%%%%%%%%%%%%%%%%%%%%%%%%%%%%%%%%%%%%%%%%%%%%%%%%%%%%%%%%%%%%%%%%%
